\chapter{Data and Feature Stores}\label{ch:ml:data-and-feature-stores}

A one-second order-book forecast has a rank information coefficient of 0.44 in its backtest and 0.36 live. Nothing is
wrong with the model. In research its features were computed on one-second bars and joined to each decision at the end of
its second; in production they are computed from every message up to the decision. The research features had seen up to a
second of the future, and the model had learned to use it. This chapter is about making that impossible: declaring each
feature once and computing it the same way for training and for trading, storing it with the time it became known, knowing
how fresh each value is, and testing continuously that the two computations agree. Its measurements use the book's synthetic
order-book market and four planted skews; the parity test catches all four, and only one of them costs the model anything,
which is a reason to catch them all rather than guess which matter.

\section{The same feature twice: offline and online}

\begin{definition}[Feature store, offline store, online store, feature definition]\label{def:ml:data-and-feature-stores:store}
A \emph{feature store}\index{feature store} is the system that computes, stores and serves a firm's model inputs from one
set of declarations. Its \emph{offline store}\index{offline store} holds feature values for history, keyed by entity and
time, for training and backtests; its \emph{online store}\index{online store} holds the latest values, updated as data
arrive, for models in production. A \emph{feature definition}\index{feature definition} declares a feature as data (its
inputs, its aggregation, its window, its clock and its version) so that both stores compute it from the same declaration
(Sculley and co-authors, 2015, on why this matters; Hermann and Del Balso, 2017, for an early production design).
\end{definition}

The chapter declares twelve features of the order book (\cref{tab:ml:fs:skew}): the spread, the queue and depth imbalances
and the weighted mid (the last values after each message, from Book~7's streaming \texttt{firm.lobfeat} engine), order-flow
imbalance over five and thirty seconds, traded and signed volume over thirty seconds, trade and message counts over ten, and
the mid's change over five and thirty. Each is a \texttt{FeatureDef}: a name, the per-event input, an aggregation (last,
sum, count, change) and a window. The \omterm{def:ml:data-and-feature-stores:store}{offline store} computes all twelve for a whole history at a list of decision times at
once, from cumulative sums and as-of lookups (\cref{lst:ml:fs:offline}); the online engine consumes the messages one by one
and keeps running windows (\cref{lst:ml:fs:online}). On eight thirty-minute sessions of \texttt{firm.tape}, 313\,681 messages
and 7\,714 decisions (one at each trade), the two agree exactly at every decision: that is the store's contract.

\section{Point-in-time correctness in the store}

Every value in the \omterm{def:ml:data-and-feature-stores:store}{offline store} carries the time it became known, and a query for a decision at time $t$ sees only values
known at $t$: the as-of join of Book~7 (chapter~3), applied to features. For event-driven features, knowledge time is the
time of the last message the feature has read; for features from slower data (fundamentals, alternative data, chapter~15)
it is the delivery time, and the store is bitemporal. The two bugs this rules out are the ones planted here: a research
pipeline that reads one message past the decision (an off-by-one in a join), and one that computes features on a coarser
clock and attaches to each decision the value at the end of its bar.

\section{Freshness and materialisation}

\begin{definition}[Materialisation, feature freshness]\label{def:ml:data-and-feature-stores:fresh}
\emph{Materialisation}\index{materialisation} is the job that computes feature values from their definitions and writes them
to a store, in batch for the \omterm{def:ml:data-and-feature-stores:store}{offline store} or continuously for the online one. \emph{Feature freshness}\index{feature
freshness} is the age of a served value: the time since the newest input it reflects, which a production system monitors
and bounds.
\end{definition}

Freshness is where production differs from research even when the code is shared: an input feed that arrives late makes every
online value that depends on it stale. The fourth planted skew delays trade prints by 200~milliseconds relative to the book
updates, as a separate trade feed might; the three trade-count and volume features then differ from research at 92 to 94\% of
decisions, by 7 to 17\% of their standard deviation, while the book features are untouched.

\section{Testing for skew}

\begin{definition}[Training--serving skew]\label{def:ml:data-and-feature-stores:skew}
\emph{Training--serving skew}\index{training--serving skew} is any difference between the feature values a model was trained
on and those it receives in production for the same moment: different code, different clocks, different inputs, different
arithmetic or different data timing.
\end{definition}

The parity test (\cref{lst:ml:fs:parity}) recomputes production's values offline for a sample of recent decisions and compares
them feature by feature. \Cref{tab:ml:fs:skew} shows each planted skew's footprint, and \cref{tab:ml:fs:parity} what the test
and the model see. A check of five random decisions per session, run five times on each of the eight sessions, flags every
skew every time: each changes most decisions. The tolerance is a design choice: at $10^{-9}$ the test also flags production's
float32 storage (a change of at most $10^{-7}$, harmless to the model), and at $10^{-6}$ it stops flagging it and still
flags the other three at 96 to 100\% of decisions.

\begin{table}[htbp]
\centering\small
\begin{tabular}{@{}lcccc@{}}
\toprule
& \multicolumn{4}{c}{RMS difference as a share of the feature's standard deviation}\\
\cmidrule(l){2-5}
feature & one event ahead & one-second bars & float32 & late trades\\
\midrule
spread & 0.38 & 0.92 & 0 & 0\\
imbalance & 0.24 & 0.73 & 0.00 & 0\\
depth imbalance & 0.11 & 0.41 & 0.00 & 0\\
weighted mid minus mid & 0.30 & 0.79 & 0.00 & 0\\
OFI 5 s & 0.03 & 0.24 & 0 & 0\\
OFI 30 s & 0.01 & 0.07 & 0 & 0\\
volume 30 s & 0.01 & 0.07 & 0 & 0.07\\
signed volume 30 s & 0.02 & 0.10 & 0 & 0.11\\
trades 10 s & 0.02 & 0.18 & 0 & 0.17\\
mid change 5 s & 0.09 & 0.39 & 0 & 0\\
mid change 30 s & 0.03 & 0.13 & 0 & 0\\
messages 10 s & 0.00 & 0.04 & 0 & 0\\
\bottomrule
\end{tabular}
\caption{Size of each planted skew, feature by feature, over 7\,714 decisions (0.00: a difference smaller than half a
hundredth; 0: none). Data: \texttt{ml\_featstore.feature\_skew}.}\label{tab:ml:fs:skew}
\end{table}

\begin{table}[htbp]
\centering\small
\begin{tabular}{@{}lccccc@{}}
\toprule
& decisions & detected & \multicolumn{2}{c}{rank IC} & live P\&L\\
\cmidrule(lr){4-5}
skew & mismatched & (5-sample check) & backtest & live & (ticks)\\
\midrule
none & 0 & 0 & 0.411 & 0.411 & 0.089\\
research one event ahead & 100\% & 100\% & 0.395 & 0.418 & 0.091\\
research on one-second bars & 96\% & 100\% & 0.438 & 0.360 & 0.078\\
production float32 storage & 92\% & 100\% & 0.411 & 0.411 & 0.089\\
production trade prints 200 ms late & 100\% & 100\% & 0.411 & 0.411 & 0.089\\
\bottomrule
\end{tabular}
\caption{Parity test (tolerance $10^{-9}$) and a ridge forecast of the mid's change over the next second, trained on four
sessions with each research pipeline and tested on four others: its backtest on research features and its live result on
production features, rank IC and mean P\&L per decision of trading the forecast's sign. Data:
\texttt{ml\_featstore.parity\_table}, \texttt{model\_effects}.}\label{tab:ml:fs:parity}
\end{table}

\begin{omfigure}
\begin{tikzpicture}
\begin{axis}[width=11.5cm,height=5.8cm,ybar,bar width=10pt,ymin=0.3,ymax=0.46,ylabel={rank IC, one-second forecast},area legend,
  xtick={0,1,2,3,4},xticklabels={{none},{one event\\ahead},{one-second\\bars},{float32\\storage},{late trade\\prints}},
  xticklabel style={align=center},tick label style={font=\small},label style={font=\small},xmin=-0.5,xmax=4.5,grid=major,
  grid style={gray!25},legend columns=2,legend style={font=\footnotesize,at={(0.5,-0.32)},anchor=north,draw=none,fill=none},
  table/col sep=comma]
\addplot[fill=omThm,draw=none,bar shift=-5.5pt] table[x=i,y=backtest]{figdata/ml/24-data-and-feature-stores/ic.csv};
\addplot[fill=omDef,draw=none,bar shift=5.5pt] table[x=i,y=live]{figdata/ml/24-data-and-feature-stores/ic.csv};
\legend{backtest (research features),live (production features)}
\end{axis}
\end{tikzpicture}

\omcaption{What each skew does to the model: the rank IC its research backtest reports and the one it earns live. Data:
\texttt{ml\_featstore.model\_effects}.}\label{fig:ml:fs:ic}
\end{omfigure}

The chapter's named result is in the table's third row. Research on one-second bars inflates the backtest's IC from 0.411 to
0.438 and, because the model learned to lean on features that had seen the future, delivers 0.360 live, below what the clean
pipeline earns; the P\&L per decision falls from 0.089 ticks to 0.078. The one-event look-ahead does the reverse here: the
message after a trade is usually a quote update deep in the book, which adds noise rather than information, and the backtest
understates the live result (0.395 against 0.418). On the five-second forecast the same skews move the IC by less than 0.01:
whether a skew costs anything depends on how its error correlates with the target, which is not knowable in advance. The
store's job is to make the stores agree, not to judge which disagreements are harmless.

\begin{method}[Running features in production]\label{met:ml:fs:method}
\begin{enumerate}
\item Declare every feature once, as data, with its clock and window, and version it.
\item Compute offline and online from the same declarations; in the \omterm{def:ml:data-and-feature-stores:store}{offline store}, stamp every value with its knowledge time
and join point in time.
\item Run a parity test on a sample of recent decisions every day, with tolerances set per feature, and alert on any mismatch.
\item Monitor each online feature's freshness against a bound, and fall back or halt when inputs go stale.
\end{enumerate}
\end{method}

\section{Tutorial: same features twice}\label{tut:ml:data-and-feature-stores:tut}

\begin{tutorial}
\textbf{Goal.} Declare twelve features, compute them offline and online on eight sessions, plant four skews, and measure what
the parity test and the model see. \textbf{End state:} \cref{tab:ml:fs:skew,tab:ml:fs:parity}, \cref{fig:ml:fs:ic}.
\begin{tutsteps}
\item \textbf{The \omterm{def:ml:data-and-feature-stores:store}{offline store}: every decision at once.}
\omcode{code/firm/featstore/firm_featstore.py}{61}{79}{Vectorised point-in-time feature values.}\label{lst:ml:fs:offline}
\item \textbf{The online engine: running windows.}
\omcode{code/firm/featstore/firm_featstore.py}{108}{124}{Online values at any moment.}\label{lst:ml:fs:online}
\item \textbf{The parity test.}
\omcode{code/firm/featstore/firm_featstore.py}{143}{151}{Comparing the stores on a sample of decisions.}\label{lst:ml:fs:parity}
\item \textbf{Run} \texttt{ml\_featstore.parity\_table()}, \texttt{parity\_table(tol=1e-6)}, \texttt{model\_effects()},
\texttt{model\_effects(5.0)}, \texttt{feature\_skew()} and \texttt{fig\_featstore.py}.
\end{tutsteps}
\textbf{What to change next.} Plant a skew that affects only one decision in a thousand and find the sample size a daily check
needs to catch it within a week.
\end{tutorial}

\section{Build: the feature store}\label{bld:ml:data-and-feature-stores:featstore}

\begin{build}
\bfield{Purpose} Model inputs computed once by definition, point in time offline, fresh online, and continuously checked
against each other.
\bfield{Interface} \texttt{FeatureDef(name, input, agg, window, version)}, \texttt{events\_from\_tape(tape, levels)} (on
\texttt{firm.lobfeat}), \texttt{offline(events, defs, times, lag\_events)}, \texttt{OnlineEngine(defs)} with \texttt{on} and
\texttt{values}, \texttt{stream(events, defs, times)}, \texttt{parity(off, on, tol, sample, seed)}, \texttt{freshness(engine,
now)}.
\bfield{Rules} No feature exists outside a definition; the \omterm{def:ml:data-and-feature-stores:store}{offline store} never returns a value known after the query time;
parity runs daily; stale inputs are alarms.
\bfield{Acceptance tests} \texttt{code/firm/featstore/tests/}: offline equals online exactly on a session; each window
aggregation matches a brute-force recomputation; a one-event look-ahead is flagged; freshness grows when inputs stop.
\bfield{Stretch} A bitemporal \omterm{def:ml:data-and-feature-stores:store}{offline store} on \texttt{firm.pit} for slow features; tolerances per feature learned from the
arithmetic; C++ or Rust online engines checked against the Python reference, as Book~7 did for \texttt{firm.lobfeat}.
\end{build}

\begin{omsources}
\item D.~Sculley and co-authors, ``Hidden technical debt in machine learning systems'', \emph{Advances in Neural Information
Processing Systems} 28, 2015.
\item J.~Hermann and M.~Del Balso, ``Meet Michelangelo: Uber's machine learning platform'', Uber engineering blog, 2017.
\item N.~Polyzotis, S.~Roy, S.~E.~Whang and M.~Zinkevich, ``Data lifecycle challenges in production machine learning: a
survey'', \emph{SIGMOD Record} 47(2), 2018.
\end{omsources}

\section{Exercises}

\begin{exercise}[$\star$]\label{exo:ml:data-and-feature-stores:1}
A thirty-second volume feature is computed at a decision at 10:00:07.300. Which trades does the \omterm{def:ml:data-and-feature-stores:store}{online store} include, and which
does a research pipeline on one-second bars joined to the end of the decision's second include?
\end{exercise}

\begin{exercise}[$\star$]\label{exo:ml:data-and-feature-stores:2}
A skew changes one decision in a thousand. How many randomly sampled decisions must a check compare to catch it with
probability 95\%?
\end{exercise}

\begin{exercise}[$\star$]\label{exo:ml:data-and-feature-stores:3}
Why does float32 storage change the imbalance features but not the volume and count features?
\end{exercise}

\begin{exercise}[$\star\star$]\label{exo:ml:data-and-feature-stores:4}
Why did the one-second bars hurt the one-second forecast and not the five-second one?
\end{exercise}

\begin{exercise}[$\star\star$]\label{exo:ml:data-and-feature-stores:5}
Why did the one-event look-ahead make the backtest worse rather than better here? Would it always?
\end{exercise}

\begin{exercise}[$\star\star$]\label{exo:ml:data-and-feature-stores:6}
\emph{Find the flaw.} ``We checked that production and research features have the same mean and standard deviation every day,
so there is no \omterm{def:ml:data-and-feature-stores:skew}{training--serving skew}.''
\end{exercise}

\begin{exercise}[$\star\star\star$]\label{exo:ml:data-and-feature-stores:7}
\emph{Coding.} Retrain the one-second model on the values the production engine computes when it replays the training
sessions, instead of the one-second-bar research features, and score it live. What does it recover, and what does it cost?
\end{exercise}

\begin{exercise}[$\star\star\star$]\label{exo:ml:data-and-feature-stores:8}
Design the tolerance of the parity test for a sum over a window of float64 inputs whose online and offline computations add in
different orders.
\end{exercise}

\section{Problem: Same Features Twice}

\begin{problem}[{Weekend problem --- one definition, two computations}]\label{pb:ml:data-and-feature-stores:1}
The chapter's twelve features, eight sessions and four skews.

\medskip\noindent\textbf{Part I --- The store.}
\begin{enumerate}
\item What does a \omterm{def:ml:data-and-feature-stores:store}{feature definition} contain, and why as data?
\item How do the offline and online computations differ, and how do we know they agree?
\item What is the knowledge time of an order-book feature, and of a quarterly fundamental?
\item What is freshness, and what made it fail here?
\end{enumerate}

\medskip\noindent\textbf{Part II --- The skews.}
\begin{enumerate}[resume]
\item Describe the four skews and which features each touches.
\item Which are research bugs and which production bugs?
\item What share of decisions does each change?
\item How does the parity test's tolerance change what it flags?
\end{enumerate}

\medskip\noindent\textbf{Part III --- The model.}
\begin{enumerate}[resume]
\item What are the backtest and live ICs under each skew?
\item Why does the one-second-bar skew inflate the backtest and hurt the live result?
\item What changes at a five-second horizon?
\item What would the IC gap look like on a desk's dashboard?
\end{enumerate}

\medskip\noindent\textbf{Part IV --- The verdict.}
\begin{enumerate}[resume]
\item State the \emph{named result}: the IC and P\&L lost to the misaligned window, and the parity test's detection rate for each
planted skew.
\item Which skew is the most dangerous, and why?
\item How would you roll out a change to a feature's definition?
\item What would you do on the day the parity test fires?
\item How should a vendor's features be brought into the store?
\item What does the store owe the model's validator (chapter~21)?
\item What is the cost of a \omterm{def:ml:data-and-feature-stores:store}{feature store} for a small team, and when is it worth it?
\item In one sentence: what does a \omterm{def:ml:data-and-feature-stores:store}{feature store} guarantee?
\end{enumerate}
\end{problem}

\section{Interview questions}

\begin{interviewq}[$\star$ \iqroles{mle}]\label{iq:ml:data-and-feature-stores:1}
What is \omterm{def:ml:data-and-feature-stores:skew}{training--serving skew}? Give three causes.
\end{interviewq}

\begin{interviewq}[$\star\star$ \iqroles{mle, dev}]\label{iq:ml:data-and-feature-stores:2}
Design a \omterm{def:ml:data-and-feature-stores:store}{feature store} for order-book features used by a model that trades on every message.
\end{interviewq}

\begin{interviewq}[$\star\star$ \iqroles{researcher, mle}]\label{iq:ml:data-and-feature-stores:3}
Your model's live IC is half its backtest's. How do you find out whether the features are to blame?
\end{interviewq}

\begin{interviewq}[$\star\star$ \iqroles{mle}]\label{iq:ml:data-and-feature-stores:4}
What is a point-in-time join, and how does a \omterm{def:ml:data-and-feature-stores:store}{feature store} implement it?
\end{interviewq}

\begin{interviewq}[$\star\star$ \iqroles{mle, dev}]\label{iq:ml:data-and-feature-stores:5}
How do you monitor \omterm{def:ml:data-and-feature-stores:fresh}{feature freshness}, and what do you do when a feature goes stale?
\end{interviewq}

\begin{interviewq}[$\star\star\star$ \iqroles{mle}]\label{iq:ml:data-and-feature-stores:6}
How would you make an online and an offline implementation of a windowed sum agree bit for bit?
\end{interviewq}
